\chapter{Fondamenti teorici e stato dell'arte}
\label{chap:stato-arte}

Il capitolo presenta i fondamenti teorici e i lavori su cui la tesi è costruita. Ogni tema è trattato come un percorso: dai lavori che ne hanno definito il problema a quelli più recenti, fino al metodo che la tesi riprende. Il percorso del capitolo segue quello del lavoro, che, partito da dati fotovoltaici PVGIS (Sezione~\ref{sec:dati-pvgis}) con il rilevatore MTGFlow, è passato a STGAN per rappresentare il contesto spaziale fra luoghi diversi e ai campi ERA5 per disporre di più grandezze meteorologiche. Per ciascun lavoro sono indicati ciò che propone, ciò che ottiene e il limite che lascia aperto, che spesso è il punto da cui parte il lavoro successivo. Che cosa ne prende la tesi è detto alla fine di ogni tema e, per STGAN, nella Sezione~\ref{sec:rw-da-stgan}. I lavori su cui il rilevatore si basa direttamente sono descritti più a fondo, con le relative formule. I temi sono quattro: il rilevamento di anomalie in serie temporali, l'apprendimento automatico applicato ai dati meteorologici, il contesto spaziale su griglie e grafi, la quantificazione dell'incertezza. Una sezione descrive poi il modo in cui la tesi adatta STGAN, e l'ultima riassume la posizione della tesi rispetto ai lavori esaminati.

\section{Rilevamento di anomalie in serie multivariate e spazio-temporali}
\label{sec:rw-anomalie}

Questa sezione va dalle definizioni di anomalia ai metodi basati su reti profonde, e arriva a STGAN, la base del rilevatore. Lungo il percorso due criteri guidano la scelta del rilevatore: la tolleranza alla presenza di anomalie nei dati di addestramento, che senza etichette non si possono escludere, e la rappresentazione del contesto spaziale fra luoghi diversi. La sezione si chiude con il problema della valutazione e con i metodi semplici usati come confronto.

\subsection{Metodi classici e modelli profondi}
\label{sec:rw-anomalie-origini}

Nella rassegna di Gupta et al.~\cite{gupta2014outlier} sui dati temporali, un'anomalia è definita in generale come una deviazione dal valore atteso o previsto. Per i dati spazio-temporali la rassegna riprende una definizione più specifica: un'osservazione è anomala quando le sue grandezze differiscono in modo significativo da quelle degli oggetti vicini nello spazio e nel tempo. I metodi che la rassegna raccoglie per questi dati procedono per lo più in due passi, prima cercando le anomalie spaziali e poi verificandole rispetto ai vicini temporali, con vicinati definiti da un raggio, da una griglia o da un raggruppamento dei dati; fra le applicazioni figurano anche dati ambientali e meteorologici, come serie di precipitazione e campi di rianalisi. I metodi spazio-temporali esaminati precedono l'uso diffuso delle reti profonde e, negli esempi riportati, trattano per lo più una sola grandezza, non campi di più variabili considerati insieme.

La rassegna di Zamanzadeh Darban et al.~\cite{zamanzadehdarban2024survey} raccoglie i metodi profondi per serie temporali e li divide secondo il segnale da cui ricavano lo \emph{score}: la \emph{previsione} del valore successivo, la \emph{ricostruzione} dell'osservazione, una rappresentazione appresa, oppure una combinazione delle prime due. Distinzioni analoghe compaiono nei lavori più recenti esaminati più avanti~\cite{wu2025catch,alnegheimish2025m2ad}. Le reti generative avversarie vi figurano fra i metodi a ricostruzione; gli autori ne indicano come limite l'addestramento, che richiede di bilanciare generatore e discriminatore ed è instabile e difficile da far convergere. Nella rassegna il termine spaziale indica le dipendenze fra le variabili di una serie, non fra luoghi geografici; fra i dataset raccolti non compaiono campi meteorologici su griglia, e i lavori su dati geografici, come STGAN, sono citati soltanto fra le applicazioni urbane e al traffico.

Due lavori mostrano i due segnali principali e il loro limite comune, cioè l'assenza di relazioni fra le serie. Malhotra et al.~\cite{malhotra2016lstm} propongono un \emph{encoder--decoder} basato su LSTM (\emph{long short-term memory}), addestrato a ricostruire finestre di serie non anomale: una finestra anomala, che il modello non ha imparato a ricostruire, produce un errore più alto, e lo score è il quadrato della distanza di Mahalanobis del vettore degli errori dalla loro distribuzione sui dati non anomali. Quando sono disponibili sequenze anomale etichettate, la soglia di decisione è però scelta massimizzando l'F-score su di esse, come in quattro dei cinque dataset dell'articolo, e nei due dataset con più sensori, i dati del motore, le serie sono ridotte alla prima componente principale: negli esperimenti, quindi, il metodo non è interamente unsupervised e non modella più variabili insieme, pur essendo formulato per vettori di più variabili. Hundman et al.~\cite{hundman2018detecting} seguono la via della previsione per la telemetria di veicoli spaziali: una LSTM per canale prevede il valore successivo, e gli errori, smussati, sono confrontati con una soglia scelta senza etichette, che separa i valori più alti dal resto della distribuzione degli errori. Gli autori indicano fra i limiti le interazioni fra i canali, che un modello per canale non rappresenta, e osservano che la qualità del rilevamento dipende da quella della previsione.

Anche i modelli più complessi non risolvono questo limite. Xu et al.~\cite{xu2022anomalytransformer} introducono Anomaly Transformer, in cui l'attenzione di un Transformer è affiancata da una distribuzione a priori concentrata sui punti vicini. Il criterio è che un'osservazione anomala, essendo rara, non riesce a stabilire legami con il resto della serie e concentra l'attenzione sui punti adiacenti; lo score moltiplica l'errore di ricostruzione per una funzione della differenza fra le due distribuzioni. La soglia è fissata senza etichette, come la percentuale di score più alti su un periodo di validazione. Le variabili sono riunite in un unico vettore per istante, l'attenzione opera solo lungo il tempo e lo score è uno per istante: il metodo non rappresenta in modo esplicito relazioni fra le serie né relazioni spaziali, e non localizza l'anomalia fra le variabili. I risultati dell'articolo sono misurati con un criterio che considera rilevato un intero segmento anomalo quando ne è segnalato anche un solo punto~\cite{xu2022anomalytransformer}.



\subsection{Relazioni fra le serie e dati di addestramento contaminati}

Una linea di lavori rappresenta le relazioni fra le serie con un grafo appreso dai dati, quando queste relazioni non sono note. Deng e Hooi~\cite{deng2021gdn} propongono GDN per i sensori di impianti industriali: a ogni sensore è associato un vettore appreso, i vicini di ciascun sensore sono i più simili secondo questi vettori, e un'attenzione su grafo prevede il valore successivo di ogni sensore. L'errore di ogni sensore è normalizzato con la sua mediana e il suo intervallo interquartile, e lo score di un istante è il massimo sui sensori, perché un'anomalia può interessarne anche uno soltanto; la soglia è il massimo dello score su un periodo di validazione. Il grafo collega sensori di uno stesso impianto, non luoghi, e si assume, come implicitamente in STGAN, che i dati di addestramento siano non anomali.

Zhou et al.~\cite{zhou2023labelfree} rimuovono questa ipotesi: osservano che gran parte dei metodi presuppone dati di addestramento composti da sole osservazioni non anomale, e che quando contengono anomalie le prestazioni possono degradare sensibilmente. Il loro metodo, MTGFlow, stima la densità dell'intero insieme di addestramento e identifica come anomale le osservazioni di test a bassa densità, sotto l'ipotesi che le anomalie occupino regioni più sparse. Le dipendenze fra entità sono apprese come grafo dinamico e un \emph{normalizing flow} specifico per entità associa a ciascuna una distribuzione gaussiana parametrizzata. La soglia è ricavata dall'intervallo interquartile degli score di addestramento e lo score per entità consente di localizzare l'anomalia.

Le entità di MTGFlow sono i sensori o le variabili di un unico sistema e la valutazione usa dataset pubblici di riferimento: il grafo è appreso dai dati e descrive relazioni fra variabili, non fra siti geografici distinti. Un'estensione assegna una distribuzione comune alle entità di uno stesso gruppo. Gli autori motivano i gruppi con la posizione comune delle entità, ma in mancanza di gruppi noti li ricavano dalla forma delle serie, con KShape. Il gruppo indica soltanto una distribuzione condivisa: il modello non usa coordinate, distanze o propagazione nello spazio, e non ha una topologia geografica esplicita.

MTGFlow è stato il primo rilevatore esaminato nel lavoro, per la sua robustezza alla presenza di anomalie nei dati di addestramento: per le condizioni meteorologiche non sono disponibili etichette, quindi non è possibile garantire dati di addestramento composti da sole osservazioni non anomale. È stato applicato ai dati PVGIS del lavoro preliminare. Non rappresenta però il contesto geografico.

I due lavori più recenti esaminati imparano il comportamento non anomalo dai dati di addestramento. Wu et al.~\cite{wu2025catch} propongono CATCH, un metodo basato su ricostruzione nel dominio della frequenza. Lo spettro è suddiviso in bande, trattate come \emph{patch}, per cogliere caratteristiche a grana fine; un modulo di fusione dei canali, con un generatore di maschere e un'attenzione mascherata, seleziona per ciascuna banda le correlazioni fra canali ritenute utili e isola quelle irrilevanti. Il metodo mira a rilevare insieme anomalie puntuali e di sottosequenza ed è valutato su dataset reali e sintetici. Lo score di anomalia combina gli errori di ricostruzione nel tempo e in frequenza: il modello impara a ricostruire i pattern non anomali, e nei dataset sintetici dell'articolo le anomalie sono inserite soltanto nella serie di test. Il caso di anomalie presenti nei dati di addestramento non è trattato. I canali sono le variabili di una serie multivariata: non è prevista una struttura geografica fra serie di siti diversi. Gli autori segnalano inoltre che anomalie più lunghe della finestra osservata non vengono colte efficacemente.

Alnegheimish et al.~\cite{alnegheimish2025m2ad} propongono M$^2$AD, per serie multivariate provenienti da più sistemi. Un modello LSTM apprende il comportamento atteso dei sensori. I residui fra valori attesi e osservati sono aggregati in uno score globale mediante un modello a miscela gaussiana e una calibrazione Gamma. La soglia è stimata dagli score osservati in addestramento, e i sensori sono ordinati per contributo allo score. Il metodo assume che i dati di addestramento siano privi di anomalie e che addestramento e test abbiano andamenti simili; gli autori indicano come limite il caso in cui la distribuzione cambia, che richiede un nuovo addestramento.

CATCH e M$^2$AD sono stati presi in esame come alternative a MTGFlow e non sono stati adottati: non soddisfano nessuno dei due criteri. Non soddisfano il primo criterio, perché M$^2$AD assume dati di addestramento privi di anomalie, una condizione che senza etichette non è verificabile sui dati meteorologici, e CATCH non tratta il caso in cui ve ne siano; né il secondo, perché, come MTGFlow, non rappresentano relazioni fra luoghi geografici. Nessuno dei due è stato eseguito sui dati della tesi. Dei due criteri della sezione, quindi, il primo ha portato a MTGFlow, il secondo, che nessuno di questi metodi soddisfa, a STGAN, che offre il contesto spaziale al prezzo di presupporre dati di addestramento non anomali.

\subsection{Reti generative avversarie per il rilevamento di anomalie}
\label{sec:rw-gan}

Una rete generativa avversaria (GAN, \emph{generative adversarial network})~\cite{goodfellow2014gan} è formata da due reti, addestrate insieme con obiettivi opposti. Il generatore $G$ trasforma un rumore $z$, estratto da una distribuzione nota $p_z$, in un campione $G(z)$; i dati reali seguono la distribuzione $p_{\mathrm{data}}$ e il discriminatore $D$ restituisce la probabilità $D(x)$ che un campione $x$ provenga dai dati reali anziché dal generatore. Le due reti sono addestrate sulla stessa funzione, che il discriminatore massimizza e il generatore minimizza:
\begin{equation}
    \min_G \max_D V(D,G) = \mathbb{E}_{x\sim p_{\mathrm{data}}}\bigl[\log D(x)\bigr] + \mathbb{E}_{z\sim p_z}\bigl[\log\bigl(1-D(G(z))\bigr)\bigr].
    \label{eq:gan-minimax}
\end{equation}
Il primo termine cresce quando il discriminatore riconosce i dati reali, il secondo quando riconosce i campioni generati; il generatore è addestrato a rendere i propri campioni indistinguibili da quelli reali. Entrambe le reti sono ottimizzate con la retropropagazione. Lo stesso articolo osserva che, all'inizio dell'addestramento, il secondo termine fornisce al generatore un gradiente debole, e propone in alternativa di addestrarlo a massimizzare $\log D(G(z))$.

Schlegl et al.~\cite{schlegl2017anogan} usano per primi questo schema per il rilevamento di anomalie, con AnoGAN, su immagini di tomografia della retina. La rete è addestrata soltanto su immagini sane; poiché una GAN non fornisce la trasformazione inversa da un'immagine al rumore che la genera, per ogni immagine da valutare il rumore è cercato con un'ottimizzazione iterativa, a reti fissate. Lo score somma due termini: la differenza fra l'immagine e quella generata, e la distanza fra le rispettive caratteristiche in uno strato intermedio del discriminatore. Sui dati dell'articolo lo score completo raggiunge un'area sotto la curva ROC di 0.89; preso da solo, il termine di ricostruzione arriva a 0.88, quello del discriminatore a 0.74~\cite{schlegl2017anogan}. Li et al.~\cite{li2019madgan} estendono l'idea alle serie temporali multivariate di impianti industriali con MAD-GAN, in cui generatore e discriminatore sono LSTM e tutte le variabili sono modellate insieme, dopo una riduzione con le componenti principali. Lo score combina, con un peso, l'errore di ricostruzione, ottenuto ancora con un'inversione iterativa nello spazio latente, e l'uscita del discriminatore. L'articolo non indica come scegliere quel peso né la soglia, e i risultati migliori sono selezionati fra più configurazioni; gli autori segnalano inoltre una forte oscillazione del recall da un'epoca all'altra, che attribuiscono all'instabilità del rilevamento basato su GAN~\cite{li2019madgan}. È un rischio che vale per ogni rilevatore basato su GAN, quindi anche per quello di questa tesi, addestrato una sola volta (Sezione~\ref{sec:conc-limiti}). Nessuno dei due metodi rappresenta relazioni spaziali fra le serie. Anche i modelli di diffusione, generativi come le GAN, sono stati portati al rilevamento: D3R ricostruisce serie multivariate di server e di un impianto dopo averle sporcate con il rumore della diffusione, che fa da collo di bottiglia regolabile senza riaddestrare il modello; lo score è l'errore di ricostruzione, le relazioni modellate sono fra variabili e non fra luoghi, e la soglia è tarata con le etichette~\cite{wang2023d3r}.

Nel rilevatore della tesi lo schema è usato in forma condizionata: il generatore riceve il contesto di un'osservazione e ne produce una ricostruzione. Poiché il contesto è fatto degli istanti precedenti, la ricostruzione è una previsione a un passo, di tre ore, e il suo errore misura quanto un istante è inatteso dato il passato recente, non quanto è lontano dal clima. Un'anomalia in senso climatologico, cioè uno stato lontano dal normale stagionale, ciò che è abituale per quella cella, quel giorno dell'anno e quell'ora (Sezione~\ref{sec:ingresso-normalizzazione-era5}), e un cambiamento inatteso in tre ore sono quindi cose diverse; la tesi le confronta attraverso la distanza dal normale e l'errore di ricostruzione del generatore (Capitolo~\ref{chap:risultati}). L'uso per il rilevamento di anomalie poggia su un'ipotesi di lavoro: che la ricostruzione sia accurata nelle condizioni frequenti nei dati di addestramento e meno accurata in quelle poco rappresentate. È un'ipotesi e non una proprietà garantita dall'addestramento. Se vale, se ne ricavano due segnali: l'errore fra l'osservazione e la sua ricostruzione, e il giudizio del discriminatore sull'osservazione. Già in AnoGAN, preso da solo, il termine del discriminatore rileva meno di quello di ricostruzione. Geiger et al.~\cite{geiger2020tadgan}, con TadGAN, osservano che l'errore di ricostruzione e l'uscita del discriminatore non si possono usare insieme così come sono: li riportano entrambi a z-score e li combinano con una media pesata o con un prodotto, e trovano che il discriminatore da solo dà il risultato meno stabile.
% Fonte: Geiger et al. 2020 (arXiv 2009.07769v3, IEEE BigData 2020), letto per intero il 16/11/2026: Sez. V-C ("The
% reconstruction errors RE(x) and Critic outputs Cx(x) cannot be directly used together as anomaly scores"; z-score di
% ciascuno, combinazione convessa eq. 10 con alpha = 0.5 o prodotto eq. 11) e Sez. VI-D ("Using Critic alone is
% unstable": F1 medio 0.29, deviazione standard 0.237; il prodotto e' preferito alla combinazione convessa).
STGAN, descritto di seguito, li combina in un unico score.

\subsection{STGAN}
\label{sec:rw-stgan}

Deng et al.~\cite{deng2022stgan} propongono una rete avversaria per il rilevamento di anomalie del traffico. Il generatore prevede la dinamica non anomala combinando tre rami: un ramo recente, basato su unità ricorrenti con convoluzione su grafo, per le dipendenze spazio-temporali di breve periodo, un ramo per il trend di lungo periodo e uno per le variabili esterne. Il discriminatore stabilisce se una sequenza è reale o generata in funzione dei vicini nello spazio e nel tempo, fornendo un criterio che varia con località e istante. Dopo l'addestramento avversario entrambe le reti agiscono da rilevatore e uno score ne combina i contributi. Il grafo deriva da una topologia fisica, cioè sensori adiacenti sulla stessa strada o regioni con confini comuni, e la valutazione misura il recall rispetto a incidenti ed eventi documentati, che gli autori stessi indicano come verità di riferimento incompleta. L'articolo non discute la presenza di anomalie nei dati di addestramento: il discriminatore è addestrato considerando non anomale le sequenze reali e anomale quelle generate, quindi l'ipotesi di dati di addestramento non anomali è implicita. La Figura~\ref{fig:schema-stgan} riassume la struttura del metodo.

\begin{figure}[htbp]
    \centering
    % Schema di STGAN: generatore a tre rami, discriminatore e composizione dello
% score. Ridisegnato a partire dalla Figura 6 di Deng et al. (2022).
\begin{tikzpicture}[font=\scriptsize,
    dato/.style={draw=black!60,rounded corners=2pt,align=center,inner sep=3pt,minimum height=0.8cm,text width=1.6cm},
    modulo/.style={draw=black!80,fill=blue!10,align=center,inner sep=3pt,minimum height=0.8cm,text width=2.05cm},
    rete/.style={draw=black!80,fill=orange!14,align=center,inner sep=4pt,text width=2.1cm},
    esito/.style={draw=black!80,fill=green!10,rounded corners=2pt,align=center,inner sep=4pt},
    freccia/.style={-{Stealth[length=4.5pt,width=4pt]},black!75},
    cornice/.style={draw=black!55,dashed,rounded corners=4pt,inner sep=6pt}]
    % ---- Generatore --------------------------------------------------------
    \node[dato] (xr) at (0,2.5) {storia recente del sottografo};
    \node[dato] (xt) at (0,1.25) {storia del nodo};
    \node[dato] (xe) at (0,0) {ora e giorno};
    \node[modulo] (mr) at (2.5,2.5) {ramo recente\\GRU su grafo};
    \node[modulo] (mt) at (2.5,1.25) {ramo del trend\\LSTM};
    \node[modulo] (me) at (2.5,0) {ramo esterno\\strato denso};
    \node[modulo,text width=1.05cm] (fu) at (4.7,1.25) {fusione};
    \draw[freccia] (xr) -- (mr);
    \draw[freccia] (xt) -- (mt);
    \draw[freccia] (xe) -- (me);
    \draw[freccia] (mr.east) -- ++(0.3,0) |- ([yshift=5pt]fu.west);
    \draw[freccia] (mt) -- (fu);
    \draw[freccia] (me.east) -- ++(0.3,0) |- ([yshift=-5pt]fu.west);
    \node[cornice,fit=(xr)(xe)(fu),label={[font=\footnotesize]above:Generatore $G_\theta$}] (gen) {};
    % ---- Sequenze e discriminatore ----------------------------------------
    \node[dato,text width=1.7cm] (sf) at (7.65,2.2) {sequenza generata $\hat{S}_{v,t}$};
    \node[dato,text width=1.7cm] (sr) at (7.65,0.3) {sequenza reale $S_{v,t}$};
    \node[rete] (d) at (10.5,1.25) {GRU su grafo sulla storia e convoluzione su grafo sul dato corrente};
    \draw[freccia] (fu.east) -- ++(0.5,0) |- node[pos=0.78,above,inner sep=2pt] {$\hat{X}_{v,t}$} (sf.west);
    \draw[freccia] (sf.east) -- ++(0.3,0) |- ([yshift=5pt]d.west);
    \draw[freccia] (sr.east) -- ++(0.3,0) |- ([yshift=-5pt]d.west);
    \node[cornice,fit=(d),label={[font=\footnotesize]above:Discriminatore $D_\phi$}] (dis) {};
    % ---- Score -------------------------------------------------------------
    \node[esito] (sg) at (4.7,-2.0) {$s_G=\lVert\hat{X}_{v,t}-X_{v,t}\rVert_2$};
    \node[esito] (sd) at (10.5,-2.0) {$s_D=D_\phi(S_{v,t})-D_\phi(\hat{S}_{v,t})$};
    \node[esito] (sc) at (7.6,-3.2) {$\mathrm{score}=s_G+\lambda\,s_D$};
    \draw[freccia] (fu.south) -- (sg.north);
    \draw[freccia] (d.south) -- node[left,align=right] {output in $[0,1]$:\\$0$ reale, $1$ generata} (sd.north);
    \draw[freccia] (sg.south) |- (sc.west);
    \draw[freccia] (sd.south) |- (sc.east);
\end{tikzpicture}

    \caption{Schema di STGAN. Il generatore combina tre rami e produce la stima $\hat{X}_{v,t}$, che sostituisce l'ultimo elemento della sequenza reale; il discriminatore valuta le due sequenze. Ridisegnato a partire da Deng et al.~\cite{deng2022stgan}.}
    \label{fig:schema-stgan}
\end{figure}

Rispetto all'equazione~\eqref{eq:gan-minimax}, in STGAN il significato dell'uscita del discriminatore è invertito: le sequenze reali sono trattate come non anomale e quelle generate come anomale, e il discriminatore è spinto verso $1$ sulle sequenze generate e verso $0$ su quelle reali. Per il discriminatore l'obiettivo è lo stesso a meno di questo scambio, ma la sua uscita cresce con il sospetto di anomalia e può essere usata per costruire lo score. Il termine avversario dell'obiettivo del generatore, nell'equazione~\eqref{eq:stgan-loss-g} più sotto, è la forma alternativa dell'obiettivo del generatore proposta da Goodfellow et al.~\cite{goodfellow2014gan}, scritta in questa convenzione, a cui si aggiunge l'errore di previsione.

Sia $x_{v,t}$ l'osservazione del nodo $v$ all'istante $t$ e $X_{v,t}$ l'insieme delle osservazioni nel sottografo formato da $v$ e dai suoi vicini. Sia poi $S_{v,t}$ la sequenza reale dei valori del sottografo che termina con $X_{v,t}$, e $\hat{S}_{v,t}$ la stessa sequenza in cui l'ultimo elemento è sostituito dalla stima $\hat{X}_{v,t}=G_{\theta}(v,t)$ del generatore. Gli obiettivi di addestramento dell'articolo sono
\begin{align}
    \mathcal{L}_G(\theta) &= \sum_{s} -\log\bigl(1-D_{\phi}(\hat{S}_{v,t})\bigr) + \lambda_G\,\lVert G_{\theta}(v,t)-X_{v,t}\rVert_2, \label{eq:stgan-loss-g}\\
    \mathcal{L}_D(\phi) &= \sum_{s} -\log\bigl(D_{\phi}(\hat{S}_{v,t})\bigr) - \log\bigl(1-D_{\phi}(S_{v,t})\bigr), \label{eq:stgan-loss-d}
\end{align}
dove la somma è sui punti $s=\langle v,t,x_{v,t}\rangle$ e $\lambda_G$ bilancia l'errore di previsione e il termine avversario. Lo score combina le due reti:
\begin{align}
    s_D(v,t) &= D_{\phi}(S_{v,t}) - D_{\phi}(\hat{S}_{v,t}), \nonumber\\
    s_G(v,t) &= \lVert G_{\theta}(v,t)-X_{v,t}\rVert_2, \nonumber\\
    \mathrm{score}(v,t) &= s_G(v,t) + \lambda\, s_D(v,t). \label{eq:stgan-score}
\end{align}
La componente del discriminatore non è la sua uscita sulla sequenza reale ma la differenza rispetto alla sequenza generata: gli autori la introducono per ridurre l'effetto dell'imprevedibilità dei dati reali. Negli esperimenti dell'articolo $\lambda_G=500$, le due componenti sono normalizzate e $\lambda=1$.

Il generatore di STGAN stima i valori dell'istante corrente a partire dagli istanti precedenti, e l'articolo chiama \emph{errore di previsione} la differenza fra stima e osservazione. Nella tesi il generatore ricostruisce allo stesso modo il campo di un istante dalla storia recente, e la stessa differenza è chiamata \emph{errore di ricostruzione del generatore}: i due termini indicano la stessa grandezza. Nella classificazione della rassegna di Zamanzadeh Darban et al. lo score rientrerebbe quindi fra quelli basati sulla previsione, anche se le reti avversarie vi figurano fra i metodi a ricostruzione.

\subsection{Valutazione dei rilevatori di anomalie}
\label{sec:rw-valutazione}

I risultati dei lavori precedenti sono misurati su dataset etichettati, e la loro affidabilità dipende da quella delle etichette. Wu e Keogh~\cite{wu2021flawed} esaminano i dataset più usati per le serie temporali univariate e vi riscontrano quattro difetti ricorrenti. Il primo è che molti problemi si risolvono con una sola riga di codice basata su operazioni elementari, come media, deviazione standard e differenze successive: una riga di questo tipo risolve 316 delle 367 serie di uno dei dataset. Gli altri sono una densità di anomalie poco realistica, etichette errate o mancanti, e anomalie concentrate verso la fine delle serie. Gli autori ne concludono che parte dei progressi riportati può essere apparente, e raccomandano di confrontare ogni metodo con soluzioni banali e di esaminare sempre visivamente dati e score. L'analisi riguarda serie univariate e non tratta campi multivariati su griglia; le sue indicazioni valgono però anche in assenza di etichette, dove il confronto con metodi semplici e l'esame diretto di mappe e serie degli score restano gli strumenti disponibili.

Liu e Paparrizos~\cite{liu2024elephant} riprendono il problema con TSB-AD, un insieme di 1070 serie, univariate e multivariate, ricontrollate a mano, su cui confrontano quaranta rilevatori con gli stessi dati e la stessa procedura di scelta dei parametri. Mostrano che il criterio che considera rilevato un intero segmento quando ne è segnalato un solo punto premia anche uno score casuale, e adottano misure che non dipendono da una soglia. Con queste misure i metodi non neurali sono spesso fra i migliori, con un'analisi per componenti principali in testa fra le serie univariate; fra le multivariate prevalgono reti semplici, convoluzionali o ricorrenti, e lo stesso metodo per componenti principali, mentre Anomaly Transformer risulta ultimo in entrambi i gruppi. TSB-AD comprende anche cinque modelli di fondazione: quattro pre-addestrati su grandi raccolte di serie a prevedere o a ricostruire, come MOMENT, Chronos e TimesFM, e uno, OFA, adattato da un modello linguistico; gli autori li trovano promettenti soprattutto sulle anomalie puntuali e avvertono del rischio che le serie di valutazione siano già state viste nel pre-addestramento~\cite{liu2024elephant}. Anche TSB-AD richiede etichette, non contiene campi su griglia e non studia dati di addestramento contaminati da anomalie; ne resta l'indicazione che un rilevatore appreso va sempre confrontato con metodi semplici, perché la sua superiorità non è scontata. Zhu et al.\ lo mostrano per cinque famiglie di modelli di fondazione, MOMENT, Chronos, TimesFM, Time-MoE e TSPulse: con lo score dato dall'errore di ricostruzione o di previsione, non si distinguono in modo significativo da due baseline di una riga, la varianza su una finestra mobile per i modelli che ricostruiscono e, ispirata alle soluzioni di una riga di Wu e Keogh~\cite{wu2021flawed}, la differenza al quadrato fra un valore e la media dei precedenti per quelli che prevedono~\cite{zhu2026oneliners}. Con un solo istante precedente la seconda è il quadrato dell'errore della persistenza, l'idea della variazione in tre ore usata qui come baseline (Sezione~\ref{sec:exp-baseline}). La maggior parte di questi modelli tratta inoltre in modo nativo solo serie univariate~\cite{zhu2026oneliners}.

Per questi motivi le segnalazioni del rilevatore sono confrontate con sei modelli shallow, scelti fra i metodi non neurali valutati in TSB-AD, cioè l'analisi per componenti principali (PCA), CBLOF (\emph{cluster-based local outlier factor}), MCD (\emph{minimum covariance determinant}), HBOS (\emph{histogram-based outlier score}), COPOD (\emph{copula-based outlier detection}) e Isolation Forest, e con due metodi statistici semplici, la distanza dal normale, cioè quanto un valore si allontana dal normale stagionale della sua cella, del giorno e dell'ora, e la variazione in tre ore, cioè quanto è cambiato rispetto a tre ore prima: in tutto otto baseline (Sezione~\ref{sec:exp-baseline}). Senza etichette, un termine di paragone interno può venire dall'accordo fra i metodi. Zimek et al.~\cite{zimek2014ensembles} osservano che ogni rilevatore coglie solo alcuni aspetti di che cosa sia un'anomalia e propongono di unirne i giudizi in un consenso, come il voto a maggioranza di una giuria: è utile se i rilevatori sono abbastanza affidabili e decidono in modo indipendente, mentre fra rilevatori simili il risultato combinato tende a ripetere i loro errori comuni.
% Fonte: Zimek, Campello, Sander, SIGKDD Explorations 15(1), pp. 11-22, letto per intero il 16/11/2026: p. 11 ("every
% model ... fits only to some aspects of the 'whole truth'", "producing a consensus of judgements", "majority vote of a
% jury", "as long as the judgements are, overall, somewhat reliable and every member decides independently") e p. 14
% (con risultati simili "the combined result gets rather attracted towards the majority of the single results").
% Tolto il 1 novembre 2026: i test spaziali non sono nel confronto (Chen et al. lasciato fuori per il
% tempo di calcolo, "per adesso"; Shekhar et al. tratta un attributo alla volta). Da riprendere se tornano:
% Per il contesto spaziale esistono metodi altrettanto semplici, che confrontano il valore di un luogo con quelli dei luoghi vicini senza un modello appreso. Shekhar et al.~\cite{shekhar2001graphoutliers} definiscono così gli outlier spaziali su un grafo, con un'applicazione ai sensori di traffico di Minneapolis--Saint Paul. Per ogni nodo calcolano la differenza fra il valore di un attributo e la media dei suoi vicini, la standardizzano con la media e la deviazione standard di queste differenze su tutti i nodi e segnalano i nodi in cui supera una soglia. Il vicinato può essere spaziale, temporale, cioè lo stesso nodo negli istanti adiacenti, oppure spazio-temporale, e il test riguarda un solo attributo. Chen et al.~\cite{chen2008spatialoutliers} osservano che i test di questo tipo trattano un attributo alla volta e che un vicino con valori molto estremi può far apparire anomali anche i luoghi normali attorno. Propongono quindi di confrontare ogni luogo con la mediana dei vicini e, quando gli attributi sono più di uno, di misurare il vettore delle differenze con una distanza di Mahalanobis a stime robuste, quelle del determinante di covarianza minimo; il metodo è provato su dati del virus del Nilo occidentale. Nessuno dei due metodi apprende la dinamica dei dati né è pensato per campi continui su griglia.

\section{Apprendimento automatico e dati meteorologici}
\label{sec:rw-meteo}

Questa sezione descrive come l'apprendimento automatico è stato applicato ai campi meteorologici, prima per prevederli e poi per riconoscerne gli eventi estremi, e mostra che il rilevamento unsupervised di anomalie su tutti gli stati di una rianalisi non è stato affrontato nei lavori esaminati.

\subsection{Previsione numerica e modelli basati sui dati}
\label{sec:rw-nwp}

La previsione meteorologica numerica tratta l'evoluzione dell'atmosfera come un problema ai valori iniziali: le equazioni che la governano sono integrate a partire dallo stato osservato, e i processi di scala inferiore alla griglia, come convezione e nubi, sono rappresentati da parametrizzazioni, e lo stesso vale per la radiazione. Bauer et al.~\cite{bauer2015quiet} ne ricostruiscono i progressi, che attribuiscono soprattutto a parametrizzazioni migliori, alla previsione d'insieme e all'assimilazione dei dati, e stimano che l'accuratezza fra tre e dieci giorni sia cresciuta di circa un giorno per decennio. Fra le sfide aperte indicano la verifica delle code della distribuzione, per cui servono misure adatte a statistiche scarse. Le rianalisi come ERA5~\cite{hersbach2020era5} nascono dallo stesso apparato: combinano un modello di previsione con le osservazioni disponibili per stimare lo stato dell'atmosfera in ogni istante. Le osservazioni disponibili cambiano negli anni, e l'interazione fra gli errori sistematici del modello e un sistema di osservazione che evolve può influire sulla coerenza nel tempo del prodotto~\cite{hersbach2020era5}.
% Fonte: Hersbach et al. 2020, Sez. 10.4: "The interaction between model bias and an evolving observing system ... is a
% concern in climate reanalysis, since it can affect the temporal consistency and accuracy of the product."

Rasp et al.~\cite{rasp2020weatherbench} rendono ERA5 il riferimento comune per i modelli di previsione basati sui dati, con WeatherBench: un insieme di campi ERA5 a risoluzione ridotta, periodi di addestramento, validazione e test contigui nel tempo e confronti con persistenza, climatologia e previsione numerica. Nella discussione gli autori osservano che gli eventi estremi, essendo rari, contano poco nelle metriche medie come l'errore quadratico, che le reti estrapolano male a stati diversi da quelli visti in addestramento e che una rete addestrata con l'errore quadratico produce campi troppo regolari; propongono le GAN come possibile rimedio a quest'ultimo problema. Lam et al.~\cite{lam2023graphcast} portano l'approccio alla griglia di $0.25^\circ$ con cui ERA5 è distribuito, con GraphCast, una rete su grafo addestrata sui campi ERA5 dal 1979 al 2017, che supera la previsione deterministica operativa sul 90.3\% delle combinazioni di variabile, livello e orizzonte esaminate. Gli autori riconoscono che l'addestramento con l'errore quadratico porta il modello a esprimere l'incertezza sfumando i campi, e che il modello migliora quando è riaddestrato su anni più recenti, per ragioni non del tutto chiarite. Per valutare gli eventi estremi definiscono come tali i valori oltre un percentile di una climatologia per luogo, ora del giorno e mese. Fra i modelli basati sui dati sono comparsi anche modelli di diffusione e di fondazione: GenCast è un modello di diffusione condizionale addestrato su ERA5 dal 1979 al 2018, che genera insiemi di previsioni globali a quindici giorni~\cite{price2025gencast}; Aurora è un modello di fondazione da 1.3 miliardi di parametri, pre-addestrato su oltre un milione di ore di dati geofisici e adattato a più compiti di previsione~\cite{bodnar2025aurora}. Sono previsori, non rilevatori: gli estremi vi compaiono come obiettivo di previsione, e il loro errore di previsione non è stato usato come score, come qui quello del generatore. Gli autori di GenCast osservano che un modello di diffusione costa più di uno deterministico equivalente, perché ogni passo di previsione richiede più valutazioni della rete~\cite{price2025gencast}.

Questi lavori usano ERA5 per prevedere; la tesi lo usa per rilevare anomalie meteorologiche. Alcune loro osservazioni valgono anche per questo impiego: un modello addestrato su anni passati con un errore medio rappresenta meglio le condizioni frequenti, tende a smussare quelle rare, e il suo comportamento può cambiare quando i dati si allontanano dal periodo di addestramento.

\subsection{Eventi estremi e anomalie nei campi meteorologici}

I primi lavori che applicano reti profonde agli eventi meteorologici estremi lo fanno in modo supervised. Liu et al.~\cite{liu2016extremeweather} classificano con reti convoluzionali poco profonde ritagli di campi centrati su cicloni tropicali, fiumi atmosferici e fronti, con etichette ottenute da criteri a soglia e da esperti, e riportano accuratezze di test fra l'89\% e il 99\%; il metodo classifica ritagli già centrati e non localizza gli eventi. Gli autori indicano nella scarsità di dati etichettati un limite, e la necessità di approcci unsupervised. Racah et al.~\cite{racah2017extremeweather} passano alla localizzazione con ExtremeWeather, un insieme di simulazioni climatiche in cui i riquadri degli eventi sono ottenuti da criteri a soglia: un encoder--decoder convoluzionale nello spazio e nel tempo unisce una loss supervised sui riquadri a una di ricostruzione su tutti i campi, così che contribuiscano anche quelli senza etichette. Gli autori osservano che etichette prodotte da criteri a soglia sono incomplete, e che un metodo supervised addestrato su di esse può al più imitarli; lasciano esplicitamente aperto il rilevamento di anomalie.

Prabhat et al.~\cite{prabhat2021climatenet} affrontano il problema delle etichette con ClimateNet, un insieme di campi annotati a mano da esperti. La campagna ha coinvolto circa ottanta esperti e ha prodotto 459 immagini etichettate; l'accordo fra due esperti sulla stessa immagine, misurato come sovrapposizione delle aree annotate, è in media di 0.51, e il modello di segmentazione addestrato su queste etichette ottiene un accordo analogo con gli esperti. Le etichette sono quindi costose e, sui confini degli eventi, incerte anche per gli esperti. Il modello è addestrato e valutato su simulazioni climatiche, e l'applicazione alle rianalisi è lasciata a lavori futuri.

La scarsità di etichette motiva i lavori più recenti, che applicano modelli generativi unsupervised a ERA5. Zaninelli et al.~\cite{zaninelli2026aiattribution} addestrano un autoencoder variazionale sui campi giornalieri di temperatura a due metri di un clima di riferimento, preindustriale in un modello climatico oppure i primi decenni di ERA5, e trattano come anomalo ciò che il modello non ricostruisce nei campi recenti; l'obiettivo è l'attribuzione delle ondate di calore al cambiamento climatico, e l'anomalia è quindi la distanza da un clima di riferimento, su una sola variabile. Paçal et al.~\cite{pacal2026heatwavedrivers} addestrano un autoencoder variazionale spazio-temporale su nove variabili ERA5, con addestramento sugli anni 1941--1990 e test sugli anni 2001--2022, e raggruppano nello spazio latente le configurazioni atmosferiche delle ondate di calore in Europa occidentale; i campioni sono però selezionati in anticipo con una soglia sulla temperatura, e gli autori notano che le ricostruzioni del modello smussano gli estremi.

Nelle ricerche condotte per questa tesi non è emerso un lavoro che assegni uno score di anomalia a tutti gli stati dell'atmosfera di una rianalisi, su più variabili insieme e con un modello spazio-temporale avversario: i lavori supervised richiedono etichette che sono scarse e incerte, quelli unsupervised su ERA5 riguardano una sola variabile oppure eventi già selezionati. La ricerca non è esaustiva, e l'assenza va intesa in questo senso.

\section{Contesto spaziale su griglie e grafi}
\label{sec:rw-spazio}

STGAN rappresenta il contesto spaziale con convoluzioni su grafo, che presuppongono una rete di adiacenze. Le convoluzioni su grafo si sono diffuse nella previsione del traffico: Yu et al.~\cite{yu2018stgcn} propongono STGCN, in cui blocchi formati da una convoluzione temporale, una convoluzione su grafo e una seconda convoluzione temporale prevedono la velocità del traffico sulle stazioni di una rete stradale. Il grafo è dato dalla rete stradale oppure costruito con un kernel sulla distanza fra le stazioni, con parametri fissati a mano. Gli autori motivano il grafo osservando che una convoluzione ordinaria limita il modello a dati disposti su una griglia, come immagini e video. È proprio il caso dei campi della tesi, disposti su una griglia regolare: questa sezione richiama l'operatore che sulla griglia svolge il ruolo della convoluzione su grafo, l'unità ricorrente convoluzionale.

Le unità ricorrenti a \emph{gate}, come quella introdotta da Cho et al.~\cite{cho2014gru} e nota come GRU, aggiornano uno stato nascosto mediante gate che regolano quanta parte dello stato precedente conservare e quanta sostituire. Nella forma originaria ingresso e stato sono vettori e i gate sono trasformazioni dense. Quando l'ingresso è un campo disposto su una griglia, le trasformazioni dense possono essere sostituite da convoluzioni, così che lo stato mantenga la struttura spaziale e i parametri siano condivisi fra le posizioni. L'idea è stata proposta per la LSTM nella previsione a breve termine delle precipitazioni~\cite{shi2015convlstm} e per la GRU nell'apprendimento di rappresentazioni di sequenze video~\cite{ballas2016convgru}. STGAN applica la stessa idea a un grafo: nel suo ramo recente i prodotti matriciali della GRU sono sostituiti da convoluzioni su grafo~\cite{deng2022stgan}.

Sia $X_s$ l'ingresso all'istante $s$ e $H_{s-1}$ lo stato precedente, entrambi disposti sulla griglia. Uno strato ricorrente convoluzionale a gate calcola
\begin{align}
    R_s &= \sigma\!\left(W_r * [X_s,H_{s-1}]+b_r\right), \label{eq:convgru-reset}\\
    U_s &= \sigma\!\left(W_u * [X_s,H_{s-1}]+b_u\right),\\
    \widetilde{H}_s &= \tanh\!\left(W_c * [X_s,R_s\odot H_{s-1}]+b_c\right),\\
    H_s &= U_s\odot H_{s-1}+(1-U_s)\odot\widetilde{H}_s, \label{eq:convgru-stato}
\end{align}
dove $W_r$, $W_u$, $W_c$ sono i kernel delle convoluzioni e $b_r$, $b_u$, $b_c$ i termini costanti, appresi durante l'addestramento, $\sigma$ è la funzione sigmoidale, $*$ una convoluzione spaziale, $\odot$ il prodotto elemento per elemento e $[\cdot,\cdot]$ la concatenazione sui canali. Il gate di reset $R_s$ regola quanto dello stato precedente entra nello stato candidato $\widetilde{H}_s$; il gate di aggiornamento $U_s$ regola quanto dello stato precedente viene conservato. La forma è quella dell'unità su grafo di STGAN, con la convoluzione su grafo sostituita da una convoluzione spaziale.

Con un \emph{kernel} di lato $k$ e \emph{padding} $\lfloor k/2\rfloor$, ogni convoluzione vede le celle entro $\lfloor k/2\rfloor$ posizioni. Il gate di reset e lo stato candidato applicano due convoluzioni in sequenza, quindi a ogni istante lo stato precedente si propaga fino a $2\lfloor k/2\rfloor$ posizioni: l'area che influenza una cella cresce con il numero di istanti e di strati. La tesi impiega questo operatore nel generatore e nel discriminatore di STGAN (Sezioni~\ref{sec:stgan-convgru} e~\ref{sec:stgan-discriminatore}).

\section{Quantificazione dell'incertezza}
\label{sec:rw-sde}

Lo score di STGAN è un numero unico per osservazione, senza una misura della sua variabilità. Nella rassegna di Abdar et al.~\cite{abdar2021review} l'incertezza di una previsione è distinta in due tipi: quella aleatoria, propria dei dati e non riducibile, e quella epistemica, dovuta alla conoscenza limitata del modello e ai dati insufficienti. Gli autori indicano le approssimazioni bayesiane e gli insiemi di modelli come i due metodi più usati, e il riconoscimento di dati lontani da quelli di addestramento come un problema ancora aperto; è l'incertezza epistemica quella che ci si attende cresca per stati poco rappresentati nei dati, come le condizioni meteorologiche rare per un modello addestrato su anni passati. Questa sezione presenta il \emph{dropout} Monte Carlo, applicato al rilevatore, e le principali alternative.

Gal e Ghahramani~\cite{gal2016dropout} interpretano l'addestramento con dropout, che a ogni passaggio spegne a caso una parte dei neuroni, come un'inferenza bayesiana approssimata. Mantenendo il dropout attivo in previsione e ripetendo più volte il passaggio in avanti si ottiene un insieme di uscite, la cui dispersione fornisce una stima dell'incertezza del modello con un costo aggiuntivo contenuto. L'equivalenza è ricavata per una rete con il dropout prima di ogni strato con parametri; nell'esperimento con una rete convoluzionale gli autori lo applicano, come d'uso in queste reti, soltanto prima dell'ultimo strato completamente connesso. Dati $J$ passaggi in avanti $\hat{y}_j(x)$ ottenuti con maschere di dropout diverse, la media predittiva è stimata da
% Gal e Ghahramani 2016 letto il 10 novembre 2026 (Desktop/1506.02142v6.pdf): 'dropout applied before every weight layer' (Sez. 1); LeNet con 'dropout applied before the last fully connected inner-product layer (the usual way dropout is used in convnets)' (Sez. 5).
\begin{equation}
    \mathbb{E}[y\mid x]\approx\frac{1}{J}\sum_{j=1}^{J}\hat{y}_j(x),
    \label{eq:mcdropout-media}
\end{equation}
e la varianza predittiva dalla varianza campionaria dei $J$ passaggi, a cui l'articolo somma l'inverso della precisione del modello.

Nella tesi il dropout Monte Carlo è applicato alle due reti del rilevatore, il generatore e il discriminatore di STGAN, non per stimare l'incertezza di una previsione ma per campionare lo score di anomalia. Durante il calcolo degli score i dropout restano attivi e il calcolo è ripetuto con maschere di dropout diverse, le \emph{estrazioni}: per ogni valore sono riportate la media dello score sulle estrazioni, come nell'equazione~\eqref{eq:mcdropout-media}, e la sua deviazione standard, la \emph{dispersione}, come descritto nella Sezione~\ref{sec:score-dropout-montecarlo}. In questo impiego la dispersione misura la variabilità dello score rispetto alle maschere di dropout e non costituisce una stima calibrata dell'incertezza; il termine di precisione previsto da Gal e Ghahramani non viene aggiunto.

Le alternative principali richiedono di cambiare l'addestramento. Lakshminarayanan et al.~\cite{lakshminarayanan2017ensembles} propongono il Deep Ensemble, un'alternativa non bayesiana: più reti vengono addestrate in modo indipendente a partire da inizializzazioni diverse, e le loro previsioni sono combinate come una miscela a pesi uguali. Il metodo è semplice da realizzare, produce un'incertezza maggiore sugli ingressi lontani dai dati di addestramento, ma richiede di addestrare una rete per ogni membro dell'insieme, mentre il dropout Monte Carlo ottiene la dispersione da un solo addestramento. Una terza famiglia di metodi ottiene la dispersione rendendo stocastica la dinamica interna della rete, modellando l'evoluzione dello stato nascosto come un'equazione differenziale stocastica, in cui una rete di drift ne descrive l'andamento e una rete di diffusione scala un termine casuale: SDE-Net~\cite{kong2020sdenet} e il suo adattamento SDE U-Net~\cite{monaco2025sdeunet}. Il rilevatore non li adotta: aggiungere a STGAN una rete di diffusione e una dinamica stocastica avrebbe allontanato molto il rilevatore dall'architettura dell'articolo, mentre il dropout Monte Carlo richiede soltanto di aggiungere strati di dropout alle reti esistenti. L'applicazione di questi metodi al rilevatore è indicata fra gli sviluppi futuri.

\section{Adattamento di STGAN ai campi meteorologici}
\label{sec:rw-da-stgan}

STGAN è nato per il traffico: il grafo deriva da una rete fisica, lo score combina due componenti e la valutazione misura il recall su incidenti documentati (Sezione~\ref{sec:rw-stgan}). La tesi ne riprende l'architettura a tre rami e gli obiettivi avversari, nella forma del codice pubblicato dagli autori (Sezione~\ref{sec:stgan-obiettivi}), e la porta su campi meteorologici privi di etichette con le quattro modifiche elencate nel Capitolo~\ref{chap:introduzione} e una scelta di obiettivo, descritte nei Capitoli~\ref{chap:metodologia} e~\ref{chap:esperimenti}.
\begin{itemize}
    \item \textbf{Griglia e campo intero.} Sulla griglia regolare le posizioni relative dei vicini sono definite: la convoluzione su grafo del ramo recente diventa una ConvGRU, che elabora in una sola passata il campo intero di un istante, mentre il discriminatore giudica la patch di ogni cella al posto del sottografo (Sezioni~\ref{sec:stgan-convgru} e~\ref{sec:stgan-discriminatore}). Il ramo delle variabili esterne riceve una codifica del tempo per cella ed è chiamato ramo del calendario.
    \item \textbf{Ingresso e pesi.} L'ingresso è riferito al normale stagionale di ogni cella, e l'errore di ogni variabile è pesato in proporzione inversa a quanto essa cambia abitualmente (Sezioni~\ref{sec:ingresso-normalizzazione-era5} e~\ref{sec:pesi-variabili-ricostruzione}).
    \item \textbf{Condizioni nuove rispetto all'addestramento.} STGAN tratta come non anomali i dati di addestramento e, a differenza di MTGFlow, non prevede un trattamento dei dati contaminati; per i campi meteorologici dati privi di anomalie non sono garantiti. La tesi cerca quindi le condizioni nuove rispetto agli anni di addestramento, secondo l'ipotesi di lavoro della Sezione~\ref{sec:rw-gan}: lo score misura la distanza da quanto il modello ha appreso, non la rarità di una condizione in assoluto, e sono segnalati i valori con lo score più alto (Sezione~\ref{sec:exp-protocollo}).
    \item \textbf{Score e dispersione.} Accanto allo score combinato, la tesi considera l'errore di ricostruzione del generatore riferito ai suoi valori abituali nella cella, nel giorno dell'anno e nell'ora, $G_{\mathrm{n}}$ (Sezione~\ref{sec:normalizzazione-climatologica-score}). Il dropout Monte Carlo nel generatore e nel discriminatore associa a ogni score una misura di dispersione, che l'articolo non prevede.
    \item \textbf{Validazione senza etichette.} Gli anni di validazione servono a scegliere la configurazione, con la distanza Maximum Mean Discrepancy (MMD)~\cite{gretton2012mmd} fra campi osservati e ricostruiti, e a stimare climatologia e soglia degli score; i loro valori più insoliti sono trattati a parte (Sezioni~\ref{sec:distanza-mmd} e~\ref{sec:maschere-validazione}). Un'idea simile è in Monaco et al.~\cite{monaco2025sdeunet}, che però escludono gli eventi intensi dai dati con cui addestrano il modello, mentre la tesi usa tutti gli anni di addestramento.
\end{itemize}


\section{Posizionamento del lavoro}
\label{sec:posizionamento}

\begin{table}[!t]
    \centering
    \footnotesize
    \setlength{\tabcolsep}{3pt}
    \caption{Candidati considerati per il rilevatore e percorso della tesi a confronto.}
    \label{tab:confronto-rilevatore}
    \vspace{0.6em}
    \begin{tabular}{@{}L{0.12}L{0.16}L{0.16}L{0.14}L{0.19}L{0.13}@{}}
        \toprule
        Lavoro & Segnale dello score & Anomalie nei dati di addestramento & Contesto spaziale & Dati e riferimento della valutazione & Dispersione \\
        \midrule
        STGAN \cite{deng2022stgan} & Errore di ricostruzione (di previsione nell'articolo) e discriminatore & Non trattate: dati reali considerati non anomali & Sì: grafo da topologia fisica & Traffico; recall su eventi documentati & Non prevista \\
        \addlinespace
        MTGFlow \cite{zhou2023labelfree} & Densità stimata & Ammesse per costruzione & No: relazioni fra variabili & Sensori; dataset etichettati & Non prevista \\
        \addlinespace
        CATCH \cite{wu2025catch} & Errore di ricostruzione nel tempo e in frequenza & Non trattate & No: canali di una serie & Serie multivariate; dataset etichettati, reali e sintetici & Non prevista \\
        \addlinespace
        M$^2$AD \cite{alnegheimish2025m2ad} & Residui di previsione aggregati & Escluse per ipotesi & No: sensori di più sistemi & Sensori di più sistemi; dataset etichettati e un caso industriale & Non prevista \\
        \addlinespace
        Tesi & Score del rilevatore $G_{\mathrm{n}}$, scelto dopo lo studio delle due componenti; score combinato di STGAN come confronto & Non trattate: restano nei dati di addestramento & Sì: vicini sulla griglia & Campi ERA5; nessuna etichetta: componenti dello score e baseline & Dropout Monte Carlo \\
        \bottomrule
    \end{tabular}
\end{table}
% La riga della tesi descrive il metodo, non un risultato: non indica che il
% rilevamento sia più accurato di quello dei lavori citati.

La tesi adatta un rilevatore avversario spazio-temporale, nato per il traffico, a campi meteorologici privi di etichette e di topologia fisica, lungo gli assi della Sezione~\ref{sec:rw-da-stgan}, e lo valuta senza etichette, confrontandolo con i metodi più semplici della Sezione~\ref{sec:rw-valutazione}. Il contributo è l'adattamento, la combinazione di questi elementi, il protocollo di valutazione e lo studio delle due componenti dello score di STGAN su questi dati, da cui viene la scelta dello score del rilevatore (Capitolo~\ref{chap:risultati}). Gli elementi di base sono preesistenti: la tesi non propone un nuovo algoritmo di rilevamento.
% Non scrivere che il metodo "risolve" i limiti dei lavori citati finché il
% Capitolo dei risultati non contiene le misure che lo sostengono.
% In assenza di una verità di riferimento indipendente, l'accuratezza del
% rilevamento non è misurata: dichiararlo qui o nei limiti.

La Tabella~\ref{tab:confronto-rilevatore} riassume i quattro candidati considerati per il rilevatore e il percorso della tesi, secondo i criteri usati nel capitolo. Le voci relative agli articoli riportano ciò che gli articoli dichiarano; dove un aspetto non è discusso, è indicato come non trattato. STGAN è la base del rilevatore; MTGFlow è stato eseguito nel lavoro preliminare sui dati PVGIS; CATCH e M$^2$AD sono stati valutati sugli articoli e non eseguiti.


I capitoli seguenti descrivono i dati e il rilevatore costruito su queste scelte.
